\chapter{Complex Differentiability}
\label{ch:iv01}

Complex differentiability is much more rigid than real differentiability. A single complex derivative must approximate increments from every direction while respecting multiplication by $i$, and this rigidity quickly leads to analyticity.

\section*{Learning goals}
The reader should be able to state the definitions precisely, prove the core results, choose valid contours or local expansions, and diagnose which domain and regularity hypotheses are essential.

\section*{Conceptual roadmap}
\[
\boxed{\text{local holomorphic structure}\;\longrightarrow\;\text{integral or series representation}\;\longrightarrow\;\text{global consequence}.}
\]

\section{Complex derivative}
A function is complex differentiable at $z_0$ if the difference quotient has one limit as $z\to z_0$ through arbitrary complex directions.

\section{Real-linear viewpoint}
Writing $f=u+iv$ identifies the real derivative with a $2\times2$ matrix; complex differentiability means that matrix is multiplication by one complex number.

\section{Directional consistency}
Approaching along real and imaginary directions gives necessary compatibility conditions.

\section{Holomorphic functions}
A function holomorphic on an open set is complex differentiable at every point there.

\section{Elementary examples}
Polynomials and rational functions away from poles are holomorphic; conjugation and modulus fail complex differentiability except at special points.

\section{Algebra rules}
Sums, products, quotients, and compositions obey familiar derivative rules where defined.

\section{Local linearization}
$f(z_0+h)=f(z_0)+f'(z_0)h+o(|h|)$ is the coordinate-free local approximation.

\section{Rigidity preview}
Holomorphicity will imply power-series expansions, infinite differentiability, and strong uniqueness principles.

\section{Core structural results}

\begin{theorem}[Complex-linear characterization]\label{thm:iv01-01}
A real-differentiable map $f:\mathbb C\to\mathbb C$ is complex differentiable at $z_0$ iff its real derivative commutes with multiplication by $i$.
\begin{proof}
If $Df$ is multiplication by $f'(z_0)$ it clearly commutes with $i$. Conversely any real-linear map commuting with $i$ is determined by its value at $1$ and equals complex multiplication by that value.
\end{proof}
\end{theorem}

\begin{theorem}[Product rule]\label{thm:iv01-02}
If $f,g$ are complex differentiable at $z_0$, then $(fg)'=f'g+fg'$.
\begin{proof}
Use the increment identity $f(z)g(z)-f_0g_0=(f-f_0)g+f_0(g-g_0)$ and divide by $z-z_0$, using continuity implied by differentiability.
\end{proof}
\end{theorem}

\begin{theorem}[Chain rule]\label{thm:iv01-03}
If $f$ is differentiable at $z_0$ and $g$ at $f(z_0)$, then $(g\circ f)'(z_0)=g'(f(z_0))f'(z_0)$.
\begin{proof}
Insert the complex-linear first-order expansions for $f$ and $g$ and combine their little-o remainders.
\end{proof}
\end{theorem}

\section{Worked examples}

\begin{example}[Complex derivative diagnostic]\label{ex:iv01-worked-01}
Give a rigorous worked analysis of Complex derivative. State the relevant holomorphic, contour, series, or singularity hypothesis and derive the central conclusion. A function is complex differentiable at $z_0$ if the difference quotient has one limit as $z\to z_0$ through arbitrary complex directions. The decisive step is to use the complex-analytic structure globally rather than reason from real-variable intuition alone.
\end{example}

\begin{example}[Real-linear viewpoint diagnostic]\label{ex:iv01-worked-02}
Give a rigorous worked analysis of Real-linear viewpoint. State the relevant holomorphic, contour, series, or singularity hypothesis and derive the central conclusion. Writing $f=u+iv$ identifies the real derivative with a $2\times2$ matrix; complex differentiability means that matrix is multiplication by one complex number. The decisive step is to use the complex-analytic structure globally rather than reason from real-variable intuition alone.
\end{example}

\begin{example}[Directional consistency diagnostic]\label{ex:iv01-worked-03}
Give a rigorous worked analysis of Directional consistency. State the relevant holomorphic, contour, series, or singularity hypothesis and derive the central conclusion. Approaching along real and imaginary directions gives necessary compatibility conditions. The decisive step is to use the complex-analytic structure globally rather than reason from real-variable intuition alone.
\end{example}

\begin{example}[Holomorphic functions diagnostic]\label{ex:iv01-worked-04}
Give a rigorous worked analysis of Holomorphic functions. State the relevant holomorphic, contour, series, or singularity hypothesis and derive the central conclusion. A function holomorphic on an open set is complex differentiable at every point there. The decisive step is to use the complex-analytic structure globally rather than reason from real-variable intuition alone.
\end{example}

\section{Solved dossiers}
The dossiers are canonical solved problems. The provenance ledger separately records which retained corpus rules guided each topic and which dossiers were newly designed.

\begin{problem}[Complex derivative diagnostic]\label{prob:iv01-dossier-01}
Give a rigorous worked analysis of Complex derivative. State the relevant holomorphic, contour, series, or singularity hypothesis and derive the central conclusion.
\end{problem}
\begin{solution}
A function is complex differentiable at $z_0$ if the difference quotient has one limit as $z\to z_0$ through arbitrary complex directions. The decisive step is to use the complex-analytic structure globally rather than reason from real-variable intuition alone.
\end{solution}

\begin{problem}[Real-linear viewpoint diagnostic]\label{prob:iv01-dossier-02}
Give a rigorous worked analysis of Real-linear viewpoint. State the relevant holomorphic, contour, series, or singularity hypothesis and derive the central conclusion.
\end{problem}
\begin{solution}
Writing $f=u+iv$ identifies the real derivative with a $2\times2$ matrix; complex differentiability means that matrix is multiplication by one complex number. The decisive step is to use the complex-analytic structure globally rather than reason from real-variable intuition alone.
\end{solution}

\begin{problem}[Directional consistency diagnostic]\label{prob:iv01-dossier-03}
Give a rigorous worked analysis of Directional consistency. State the relevant holomorphic, contour, series, or singularity hypothesis and derive the central conclusion.
\end{problem}
\begin{solution}
Approaching along real and imaginary directions gives necessary compatibility conditions. The decisive step is to use the complex-analytic structure globally rather than reason from real-variable intuition alone.
\end{solution}

\begin{problem}[Holomorphic functions diagnostic]\label{prob:iv01-dossier-04}
Give a rigorous worked analysis of Holomorphic functions. State the relevant holomorphic, contour, series, or singularity hypothesis and derive the central conclusion.
\end{problem}
\begin{solution}
A function holomorphic on an open set is complex differentiable at every point there. The decisive step is to use the complex-analytic structure globally rather than reason from real-variable intuition alone.
\end{solution}

\begin{problem}[Elementary examples diagnostic]\label{prob:iv01-dossier-05}
Give a rigorous worked analysis of Elementary examples. State the relevant holomorphic, contour, series, or singularity hypothesis and derive the central conclusion.
\end{problem}
\begin{solution}
Polynomials and rational functions away from poles are holomorphic; conjugation and modulus fail complex differentiability except at special points. The decisive step is to use the complex-analytic structure globally rather than reason from real-variable intuition alone.
\end{solution}

\begin{problem}[Algebra rules diagnostic]\label{prob:iv01-dossier-06}
Give a rigorous worked analysis of Algebra rules. State the relevant holomorphic, contour, series, or singularity hypothesis and derive the central conclusion.
\end{problem}
\begin{solution}
Sums, products, quotients, and compositions obey familiar derivative rules where defined. The decisive step is to use the complex-analytic structure globally rather than reason from real-variable intuition alone.
\end{solution}

\begin{problem}[Local linearization diagnostic]\label{prob:iv01-dossier-07}
Give a rigorous worked analysis of Local linearization. State the relevant holomorphic, contour, series, or singularity hypothesis and derive the central conclusion.
\end{problem}
\begin{solution}
$f(z_0+h)=f(z_0)+f'(z_0)h+o(|h|)$ is the coordinate-free local approximation. The decisive step is to use the complex-analytic structure globally rather than reason from real-variable intuition alone.
\end{solution}

\begin{problem}[Rigidity preview diagnostic]\label{prob:iv01-dossier-08}
Give a rigorous worked analysis of Rigidity preview. State the relevant holomorphic, contour, series, or singularity hypothesis and derive the central conclusion.
\end{problem}
\begin{solution}
Holomorphicity will imply power-series expansions, infinite differentiability, and strong uniqueness principles. The decisive step is to use the complex-analytic structure globally rather than reason from real-variable intuition alone.
\end{solution}

\begin{problem}[Complex-linear characterization proof dossier]\label{prob:iv01-dossier-09}
Prove the following theorem-level statement and identify the essential hypothesis: A real-differentiable map $f:\mathbb C\to\mathbb C$ is complex differentiable at $z_0$ iff its real derivative commutes with multiplication by $i$.
\end{problem}
\begin{solution}
If $Df$ is multiplication by $f'(z_0)$ it clearly commutes with $i$. Conversely any real-linear map commuting with $i$ is determined by its value at $1$ and equals complex multiplication by that value. This identifies the mechanism that makes the complex-analytic conclusion stronger than its real-variable analogue.
\end{solution}

\begin{problem}[Product rule proof dossier]\label{prob:iv01-dossier-10}
Prove the following theorem-level statement and identify the essential hypothesis: If $f,g$ are complex differentiable at $z_0$, then $(fg)'=f'g+fg'$.
\end{problem}
\begin{solution}
Use the increment identity $f(z)g(z)-f_0g_0=(f-f_0)g+f_0(g-g_0)$ and divide by $z-z_0$, using continuity implied by differentiability. This identifies the mechanism that makes the complex-analytic conclusion stronger than its real-variable analogue.
\end{solution}

\begin{problem}[Chain rule proof dossier]\label{prob:iv01-dossier-11}
Prove the following theorem-level statement and identify the essential hypothesis: If $f$ is differentiable at $z_0$ and $g$ at $f(z_0)$, then $(g\circ f)'(z_0)=g'(f(z_0))f'(z_0)$.
\end{problem}
\begin{solution}
Insert the complex-linear first-order expansions for $f$ and $g$ and combine their little-o remainders. This identifies the mechanism that makes the complex-analytic conclusion stronger than its real-variable analogue.
\end{solution}

\begin{problem}[Chapter synthesis]\label{prob:iv01-dossier-12}
Connect the definitions and principal theorems of this chapter into one reusable complex-analysis strategy.
\end{problem}
\begin{solution}
Complex differentiability is much more rigid than real differentiability. A single complex derivative must approximate increments from every direction while respecting multiplication by $i$, and this rigidity quickly leads to analyticity. The reusable strategy is to identify the analytic domain, choose the correct local expansion or contour representation, apply the structural theorem, and then audit singularities and topology.
\end{solution}

\section{Exercises with complete solutions}

\begin{exercise}\label{exr:iv01-01}
State and apply the central principle behind Complex derivative. Give one concrete consequence.
\end{exercise}
\begin{hint}
Begin from the exact definition or theorem hypothesis in the corresponding section.
\end{hint}
\begin{solution}
A function is complex differentiable at $z_0$ if the difference quotient has one limit as $z\to z_0$ through arbitrary complex directions. The same principle can then be reused in later contour, residue, conformal, or Riemann-surface arguments.
\end{solution}

\begin{exercise}\label{exr:iv01-02}
State and apply the central principle behind Real-linear viewpoint. Give one concrete consequence.
\end{exercise}
\begin{hint}
Begin from the exact definition or theorem hypothesis in the corresponding section.
\end{hint}
\begin{solution}
Writing $f=u+iv$ identifies the real derivative with a $2\times2$ matrix; complex differentiability means that matrix is multiplication by one complex number. The same principle can then be reused in later contour, residue, conformal, or Riemann-surface arguments.
\end{solution}

\begin{exercise}\label{exr:iv01-03}
State and apply the central principle behind Directional consistency. Give one concrete consequence.
\end{exercise}
\begin{hint}
Begin from the exact definition or theorem hypothesis in the corresponding section.
\end{hint}
\begin{solution}
Approaching along real and imaginary directions gives necessary compatibility conditions. The same principle can then be reused in later contour, residue, conformal, or Riemann-surface arguments.
\end{solution}

\begin{exercise}\label{exr:iv01-04}
State and apply the central principle behind Holomorphic functions. Give one concrete consequence.
\end{exercise}
\begin{hint}
Begin from the exact definition or theorem hypothesis in the corresponding section.
\end{hint}
\begin{solution}
A function holomorphic on an open set is complex differentiable at every point there. The same principle can then be reused in later contour, residue, conformal, or Riemann-surface arguments.
\end{solution}

\begin{exercise}\label{exr:iv01-05}
State and apply the central principle behind Elementary examples. Give one concrete consequence.
\end{exercise}
\begin{hint}
Begin from the exact definition or theorem hypothesis in the corresponding section.
\end{hint}
\begin{solution}
Polynomials and rational functions away from poles are holomorphic; conjugation and modulus fail complex differentiability except at special points. The same principle can then be reused in later contour, residue, conformal, or Riemann-surface arguments.
\end{solution}

\begin{exercise}\label{exr:iv01-06}
State and apply the central principle behind Algebra rules. Give one concrete consequence.
\end{exercise}
\begin{hint}
Begin from the exact definition or theorem hypothesis in the corresponding section.
\end{hint}
\begin{solution}
Sums, products, quotients, and compositions obey familiar derivative rules where defined. The same principle can then be reused in later contour, residue, conformal, or Riemann-surface arguments.
\end{solution}

\begin{exercise}\label{exr:iv01-07}
State and apply the central principle behind Local linearization. Give one concrete consequence.
\end{exercise}
\begin{hint}
Begin from the exact definition or theorem hypothesis in the corresponding section.
\end{hint}
\begin{solution}
$f(z_0+h)=f(z_0)+f'(z_0)h+o(|h|)$ is the coordinate-free local approximation. The same principle can then be reused in later contour, residue, conformal, or Riemann-surface arguments.
\end{solution}

\begin{exercise}\label{exr:iv01-08}
State and apply the central principle behind Rigidity preview. Give one concrete consequence.
\end{exercise}
\begin{hint}
Begin from the exact definition or theorem hypothesis in the corresponding section.
\end{hint}
\begin{solution}
Holomorphicity will imply power-series expansions, infinite differentiability, and strong uniqueness principles. The same principle can then be reused in later contour, residue, conformal, or Riemann-surface arguments.
\end{solution}

\section*{Chapter summary}
The chapter now forms part of the canonical complex-analysis chain from holomorphicity through residues, global continuation, and Riemann surfaces.
